\documentclass[twocolumn,9pt]{extarticle}
\usepackage{graphicx} % Required for inserting images
\usepackage{newStyle}
\usepackage{titling}
\usepackage[ruled,linesnumbered]{algorithm2e}
\usepackage{bm}
\usepackage{amsthm}
\newtheorem{theorem}{Theorem}[section]
\newtheorem{lemma}[theorem]{Lemma}
\newtheorem{deffinition}[theorem]{Definition}

\setlength{\droptitle}{-4em}

\title{Something}
\author{Eugene Francisco}
\date{July 2026}

\begin{document}

\maketitle

\section{Something}
\begin{algorithm}[t]
\SetKwInOut{Input}{Input}
\SetKwInOut{Params}{Params}
\SetKwInOut{Output}{Output}
\label{algo:test_statistic}
\caption{Test statistic $\phi$ (Kuditipudi et al.)}
\Input{string $y \in \mathcal{V}^*$, watermark key sequence $\xi \in \Xi^n$}
\Params{alignment cost $d$, block size $k$}
\Output{test statistic value $\phi(y,\xi) \in \mathbb{R}$}
\For{$i \in 1,\dots,\mathrm{len}(y)-k+1$}{
  \For{$j \in 1,\dots,n$}{
    $\mathbf{y}^i \gets \{y_{i+\ell}\}_{\ell=0}^{k-1}$,\quad
    $\bm{\xi}^j \gets \{\xi_{(j+\ell)\%n}\}_{\ell=0}^{k-1}$\;
    $\hat{d}_{i,j} \gets d(\mathbf{y}^i, \bm{\xi}^j)$\;
  }
}
\Return $\min_{i,j} \hat{d}_{i,j}$\;
\end{algorithm}
From SynthID-text, we have
\begin{align*}
	\mathrm{Score}(x) = \frac{1}{mT}\sum_{t = 1}^T\sum_{\rho = 1}^m g_\rho(x_t, r_\rho).
\end{align*}
Define
\begin{align}
	\mathrm{Score}(x, \xi ) = \frac{1}{mT}\sum_{t = 1}^T\sum_{\rho = 1}^m g_\rho(x_t, \xi_\rho)\label{eq:score_func}
\end{align}
and let
\begin{align}
	d(\cdot, \cdot) = -\mathrm{Score}(\cdot, \cdot).\label{eq:alignment_func}
\end{align}

\begin{theorem}
	Let $\phi $ be the output of Algorithm~\ref{algo:test_statistic} when ran on
	an unwatermarked string $y \in V^*$, watermark key $\xi \in \Xi^n$, alignment function
	as defined in Equation~\ref{eq:alignment_func} and chunk size $k$. Let $F_\phi$ and $f_\phi$ be the distribution and
	density of $\phi$ respectively. Then
	\begin{align*}
	F_\phi(x) &= 1 - (1 - F_\mathrm{Score}(x))^\mu\\
	f_\phi(x) &= \mu f_\mathrm{Score}(x)(1 - F_\mathrm{Score}(x))^{\mu - 1}
\end{align*}
where $\mu = (\mathrm{len}(y) - k + 1) \cdot n$ and $F_\mathrm{Score}$/
$f_\mathrm{Score}$ are the distributions and densities of $\mathrm{Score}(y, \xi)$ for $y$ unwatermarked and $\xi \sim \nu$ (in the $N = 2$ case, $\mathrm{Score}(y, \xi)\sim \mathrm{Binomial}(m\cdot T, 0.5)$).
\end{theorem}
\begin{proof}
	Observe that the distribution of $\phi$ in the theorem is the distribution of the minimum of $m$ i.i.d. random variables with density $f_{\mathrm{Score}}$ and distribution $F_{\mathrm{Score}}$; similarly for the density. It suffices to show then that $\{d(y^i, \xi^j): i\in 1, \ldots, \mathrm{len}(y) - k + 1, j\in 1, \ldots, n\}$ are i.i.d.

	To show this, recall from Equation~\ref{eq:score_func} that
	\begin{align}
		d(y^i, \xi^j) = -\frac{1}{mT}\sum_{t = 1}^k\sum_{\rho = 1}^mg_\rho(y_t^i, \xi_\rho^j).
	\end{align}
	Because of how the $g_\rho$ are chosen, we will be done if we can show, for any particular $y_{i_1}$, $\xi_{i_2}$, and $g_{i_3}$, that $g_{i_3}(y_{i_1}, \xi_{i_2})$ shows up in only one of the $\{d(y^i, \xi^j)\}$. In other words, we must show that, once $y_{i_1}$, $\xi_{i_2}$, and $g_{i_3}$ are picked, $i$ and $j$ of $d(y^i, \xi^j)$ are uniquely determined. 
\end{proof}
\TODO, proof isn't true. 





\end{document}